\chapter{Introduction}
\label{ch:introduction}

\section{Motivation}

Public unmet-care indicators support comparisons across European health systems. Eurostat EU-SILC provides country-year measures of self-reported unmet medical and dental care, including cost, distance, and waiting-list barriers and indicators defined over either the general population or people reporting a need for care \citep{eurostat_health_variables_metadata,eurostat_unmet_needs_statistics}.

These data do not automatically define one stable analytical population. Requiring complete regression covariates can move an analysis from a broad outcome-observed panel to a selected subset of countries and years. Weighting, country restrictions, and model-based treatment of missing values can alter the target further. Because incomplete-data methods depend on assumptions about selection and missingness, this movement must be documented rather than treated as routine preprocessing \citep{little_rubin_2019,seaman_white_2013}.

\section{Target-Population Drift}

This thesis uses \emph{target-population drift} as a working term for movement from an intended monitoring target to a different analytical target created by denominator choice, country-universe restriction, population weighting, covariate completeness, or missing-data modelling. The same poverty/social-exclusion coefficient may therefore describe an equal-country complete-case panel, a population-weighted resident panel, an EU-27 subset, an IPW-reweighted selected panel, or a model-imputed outcome-observed target.

The country universe is defined before model interpretation. The full Eurostat-available universe comprises national units observed in the selected Eurostat tables after EU and euro-area aggregates are excluded. EU-27 sensitivity is mandatory because the term ``European'' would otherwise conceal an important institutional restriction. Smaller country groups are reported only when coverage is adequate. Institutional groupings follow official EU, EEA/EFTA, enlargement, UK, and Switzerland sources \citep{eu_countries_2026,europarl_eea_switzerland_2025,european_commission_enlargement_2026,european_commission_uk_relations_2026}.

\section{Research Questions}

The thesis addresses five research questions:

\begin{description}
    \item[RQ1: Indicator dependence.] How do unmet-care monitoring conclusions change across medical versus dental care, population-denominator versus need-denominator indicators, and barrier-specific definitions such as cost, distance, and waiting lists?
    \item[RQ2: Target-population drift.] How much does complete-case covariate selection change the country-year population used for regression, and which variables, countries, years, and country groups drive this attrition?
    \item[RQ3: Coefficient stability.] Are social-vulnerability coefficients stable across complete-case, population-weighted, balanced-outcome, IPW, MAR-imputed, and MNAR-shifted estimands?
    \item[RQ4: Simulation validity.] Under semi-synthetic Eurostat-style missingness, when does complete-case fixed-effects estimation recover the reference estimand, and when does it produce bias, sign instability, interval undercoverage, or wrong substantive conclusions?
    \item[RQ5: Reporting protocol.] Can a reproducible target-population drift audit classify public aggregate-panel findings as stable, target-dependent, denominator-dependent, missingness-dependent, imputation-model-dependent, or non-identifiable?
\end{description}

Poverty or social exclusion is the main stress-test covariate. The primary outcome is observed for 608 country-years, whereas the baseline complete-case fixed-effects analysis retains 282 rows (Table~\ref{tab:complete_case_attrition_waterfall}). The empirical question is whether the aggregate poverty association remains interpretable across the explicitly defined targets.

\section{Drift-Stability Framework}

An estimand is the country-year quantity targeted by a specified outcome, denominator, row set, country universe, weighting rule, missing-data assumption, and model. Chapter~\ref{ch:framework} formalizes this object as an estimand tuple and defines two reporting quantities: the Target-Population Drift Index, which summarizes movement from the declared monitoring target, and the Conclusion Stability Index, which summarizes whether standardized coefficient conclusions remain stable across feasible estimands. The framework is a diagnostic reporting structure rather than a new estimator.

The sequence is target definition, drift measurement, coefficient-stability assessment, simulation stress testing, and final classification. Chapter~\ref{ch:data_indicators} defines the targets, Chapter~\ref{ch:framework} defines the drift-stability metrics, Chapter~\ref{ch:modeling_strategy_results} evaluates their empirical consequences, and Chapter~\ref{ch:discussion_conclusion} presents the reporting checklist.

\section{Contributions}

This thesis makes four contributions. First, it constructs an estimand registry and multi-outcome benchmark that record indicator denominators, country coverage, weighting options, and valid interpretations before modelling. Second, it measures complete-case attrition and selection distortion across countries, years, covariate blocks, and outcome definitions. Third, it formalizes target-population drift and conclusion stability through TDI, CSI, and rule-based classification, then evaluates coefficient stability across observed-data and imputed estimands and through semi-synthetic missingness simulation. Fourth, it provides a reproducible reporting protocol supported by extraction records, data dictionaries, output maps, tests, and scripted analysis.

The novelty is to treat sample construction as an empirical object rather than a technical footnote. Previous unmet-care research primarily examines social gradients, measurement validity, or health-system differences. This thesis instead asks whether an aggregate association retains the same interpretation when the denominator, country universe, weighting rule, analytical row set, or missing-data assumption changes. Its contribution is methodological-applied: it improves the reporting and interpretation of public aggregate health-monitoring analyses.

\section{Thesis Structure}

Chapter 2 connects healthcare-access theory with Eurostat/EU-SILC measurement, ecological inference, missing-data assumptions, and reporting discipline. Chapter 3 defines the data architecture and estimand registry. Chapter 4 formalizes the target-population drift and conclusion-stability framework. Chapter 5 compares monitoring targets before modelling. Chapter 6 presents the selection diagnostics, coefficient-stability analysis, and semi-synthetic simulation. Chapter 7 converts the evidence into a reporting protocol and final classification. Supporting prediction, equity, diagnostic, and reproducibility material is reported in the appendices.
